\capituloPregunta{¿Qué ocurre cuando dos factores actúan al mismo tiempo?}

\section{Motivación}

Un tratamiento rara vez actúa solo. En ingeniería, educación, manufactura o biología, una respuesta puede depender de la combinación de condiciones. Por eso no basta preguntar si un factor importa; también hay que preguntar si su efecto cambia cuando cambia otro factor.

\section{Caso integrado: profesor y materia}

El material del curso contiene un ejemplo balanceado con dos factores: profesor y materia. La variable respuesta es la percepción estudiantil. Cada combinación contiene observaciones, lo que permite estudiar:

\begin{itemize}
  \item efecto principal del profesor;
  \item efecto principal de la materia;
  \item interacción profesor--materia.
\end{itemize}

\begin{idea}
Una interacción aparece cuando el efecto de un factor depende del nivel del otro. En ese caso, los efectos principales pueden contar una historia incompleta.
\end{idea}

\section{Tabla de diseño}

\begin{table}[H]
\centering
\begin{tabular}{ccc}
\toprule
Combinación & Profesor & Materia\\
\midrule
1 & A & 1\\
2 & A & 2\\
3 & B & 1\\
4 & B & 2\\
\bottomrule
\end{tabular}
\end{table}

\section{Lectura antes del cálculo}

Antes de hacer ANOVA, el estudiante debe construir una gráfica de perfiles. Si las líneas son aproximadamente paralelas, la interacción puede ser pequeña. Si se cruzan o se separan claramente, la interacción debe interpretarse con cuidado.

\section{Actividad}

Proponga un ejemplo de ingeniería con dos factores. Para cada factor defina dos niveles y una variable respuesta. Después responda: ¿qué significaría una interacción en ese contexto?

\section{Conexión}

Para estudiar interacciones con mayor control, se construyen diseños factoriales.
